\chapter{The Straight Line}

\section{Angles}
\begin{enumerate}
	\item[20.] Inkscape vector image:
	
	\includegraphics[scale=0.6]{chapters/ch01/images/20}
	
	\item[21.] If the interior of the angle is smaller than the exterior then a congruent angle can be drawn inside the exterior without any overlap.
	
	It can't be done for interior.
	
	\item[22.] Two distinct angles can have two common sides: \textit{AOB (smaller)} and \textit{AOB (bigger)}
	
	\includegraphics[scale=0.6]{chapters/ch01/images/22}

	
	\item[23.] No
	
	\item[24.] Yes, arcs can have different radius but same center.
	
	No, one arc is smaller than the other.
	
%25:	
	\item[25.] Figure:
	
	\includegraphics[scale=0.8]{chapters/ch01/images/25}
	
	Since $CD$ is a straight line, we know from \S16.1:
	
	$\angle DOC = \angle DOA + \angle AOC = \ang{180}$
	
	$\Rightarrow \angle AOC = \ang{180} - \ang{25} = \ang{155}$
	
	We also know that:
	
	$\angle AOB = \angle AOD + \angle DOB = \ang{180}$
	
	$\Rightarrow \angle BOD = \ang{180} - \ang{25} = \ang{155}$
	
	Finally we get:
	
	$\angle BOC = \ang{108} - \ang{155} = \ang{25}$.
	
	\item[26.] \ang{25}, \ang{55}, \ang{100}, \ang{100}
	
	\item[27.] $19*19=361$, it's easy to show...
	
\end{enumerate}

\section{Perpendicular Lines}
\begin{enumerate}
	\item[28.] $\ang{14;24;44} + \ang{75;35;25} = \ang{90;0;9}$. Hence it is obtuse.
	
	\item[29.] $\times\times\times$ Confusing question $\times\times\times$
	
	\item[30.] 
\end{enumerate}